\langsection[subsection]{SEE}{Стратегия SEE}


\begin{langblock}{0.425}
	\begin{eng}
	Experienced riders remain aware of what is going on around them. They improve their riding strategy by using the Motorcycle Safety Foundation's (MSF) SEE strategy, a 3-step process used to make appropriate judgments, and apply them correctly in different traffic situations:
	
	\begin{itemize}
		\item \textbf{S}earch
		\item \textbf{E}valuate		
		\item \textbf{E}xecute
	\end{itemize}
	
	Let's examine each of these steps.
	\end{eng}
	\begin{rus}
		Опытные мотоциклисты всегда следят за тем, что происходит вокруг. Они совершенствуют свою стратегию вождения, используя парадигму SEE, разработанную Фондом безопасности мотоциклистов (MSF) "--- это 3-х этапный процесс, позволяющий принимать адекватные решения и правильно применять их в различных дорожных ситуациях:
		\begin{itemize}
			\item \textbf{S}earch (Поиск)
			\item \textbf{E}valuate (Оценка)
			\item \textbf{E}xecute (Действие)
		\end{itemize}
		
		Давайте рассмотрим каждый из этих этапов.
\end{rus}
\end{langblock}




\subsubsection*{Search}


\begin{langblock}{0.425}
	\begin{eng}
	Search aggressively ahead, to the sides and behind to avoid potential hazards even before they arise. How assertively you search, and how much time and space you have, can eliminate or reduce harm. Focus even more on finding potential escape routes in or around intersections, shopping areas, schools, and construction zones.
	
	Search for:
	
	\begin{description}[noitemsep]
		\item [Oncoming traffic] that may turn left in front of you.
		\item [Traffic] coming from the left and right.
		\item [Traffic] approaching from behind.
		\item [Hazardous] road conditions.
		
	\end{description}
	
	Be especially alert in areas with limited visibility. Visually busy surroundings could hide you and your motorcycle from others.
	\end{eng}
	\begin{rus}
		Внимательно просматривайте траекторию вперед, по бокам и сзади для избежания потенциальных опасностей еще до того, как они возникнут. От того, насколько тщательно вы это делаете зависит то, сколько времени и пространства у вас в запасе, чтобы предотвратить или уменьшить ущерб. Уделяйте больше внимания поиску вариантов для маневра на перекрестках, в районе торговых центров, школ и строительных площадок, а также в их окрестностях.
		
		Обратите внимание на:
		\begin{description}[noitemsep]
		\item [Встречный транспорт], который может повернуть налево прямо перед вами.
		\item [Транспорт], приближающийся слева и справа.
		\item [Транспорт], приближающийся сзади.
		\item [Опасные] дорожные условия.
		\end{description}
		
		Будьте особенно внимательны в местах с ограниченной видимостью. В окружении, где много объектов окружающие вас и ваш мотоцикл могут просто не заметить.
\end{rus}
\end{langblock}




\subsubsection*{Evaluate}


\begin{langblock}{0.425}
	\begin{eng}
	Think about how hazards can interact to create risk for you. Anticipate potential problems and have a plan to reduce risk.
	
	\begin{description}[noitemsep]
		\item [Road] and surface characteristics – Potholes, guardrails, bridges, telephone poles and trees won't move into your path but may influence your riding strategy.
		\item [Traffic control devices] – Look for traffic signals, including regulatory signs, warning signs, and pavement markings, to help you evaluate circumstances ahead.
		\item [Vehicles] and other traffic – May move into your path and increase the likelihood of a crash.
		
	\end{description}
	Think about your time and space requirements in order to maintain a margin of safety. You must leave yourself time to react if an emergency arises.
	\end{eng}
	\begin{rus}
		Подумайте о том, как различные опасности могут взаимодействовать друг с другом, создавая для вас риски. Продумайте возможные проблемы и план действий.
		
		\begin{description}[noitemsep]
			\item [Качество дороги] и дорожного покрытия "--- выбоины, ограждения, мосты, телефонные столбы и деревья сами по себе не возникают на вашем пути, но могут повлиять на вашу стратегию.
			\item [Средства организации] дорожного движения "--- отмечайте заранее светофоры, регулирующие и предупреждающие знаки, а также дорожную разметку, чтобы полно и заранее оценить обстановку.
			\item [Транспортные средства] и другие участники дорожного движения "--- могут выехать перед вами и увеличить вероятность аварии.
			
		\end{description}
		
		Подумайте о том, что в случае возникновения чрезвычайной ситуации, вам нужны время и дистанция для маневра. Добавьте время на реакцию.
\end{rus}
\end{langblock}




\subsubsection*{Execute}


\begin{langblock}{0.425}
	\begin{eng}
	
	Carry out your decision.
	
	To create more space and minimize harm from any hazard:
	
	\begin{itemize}[noitemsep]
		\item Communicate your presence with lights and/or horn.
		\item Adjust your speed by accelerating, stopping or slowing.
		\item Adjust your position and/or direction.
	\end{itemize}
	
	Apply the old adage one step at a time to handle two or more hazards. Adjust speed to permit two hazards to separate. Then deal with them one at a time as single hazards.
	
	Decision making becomes more complex with three or more hazards.  Evaluate the consequences of each and give equal distance to the hazards.
	
	In potential high risk areas, such as intersections, shopping areas, school and construction zones, cover the clutch and both brakes to reduce the time you need to react.
	\end{eng}
	\begin{rus}
		Действуйте в соответствии со своим планом.
		
		Не сужайте свои варианты, сведите к минимуму ущерб от любой опасности:
		
		\begin{itemize}[noitemsep]
			\item Сигнализируйте о своем присутствии с помощью световых приборов и/или клаксона.
			\item Регулируйте скорость, ускоряясь, останавливаясь или замедляясь.
			\item Корректируйте свое положение на полосе и/или направление движения.
		\end{itemize}
		
		Используйте старую поговорку «шаг за шагом», чтобы справиться с двумя или более проблемами. Регулируйте скорость так, чтобы две опасности разделились. Затем решайте их по очереди, как отдельные проблемы.
		При наличии трех и более опасностей принятие решений усложняется. Оцените последствия каждой и равномерно разнесите их по времени.
		
		В потенциально опасных зонах, таких как перекрестки, торговые районы, школьные и строительные зоны, чтобы сократить время реакции "--- держите пальцы на рычагах сцепления и переднего тормоза, а ногу "--- на педали заднего тормоза\footnote{В отличие от автомобилей, при экстренном торможении на мотоцикле сцепление нужно выжимать полностью и сразу.}.
\end{rus}
\end{langblock}






























